% =====================================================================
\section{Synthetic Benchmark Results}
\label{sec:results-synth}
% =====================================================================

We organise the synthetic results along two axes: spatial and temporal
complexity.  We report attribution error EC1 against oracle Shapley values
(Sec.~\ref{sec:results-error}), the 2$\times$2 spatiotemporal failure-mode
matrix (Sec.~\ref{sec:results-2x2}), and the faithfulness and ranking
metrics that transfer to the real-data setting
(Sec.~\ref{sec:results-faithfulness}).

\subsection{Attribution Error vs.\ Oracle Shapley Values}
\label{sec:results-error}

Table~\ref{tab:synth_ec1} reports EC1 across five pillar datasets spanning
the 2$\times$2 spatial $\times$ temporal design and the heavy-tailed Burr
regime, averaged over three classifier architectures (MLP, CNN,
Transformer).  Four patterns are consistent across architectures:
\textbf{(i)} \emph{Learned conditional imputers recover oracle Shapley
values to within Monte Carlo noise.}  KS--VAEAC achieves EC1 $\le 0.007$
on every pillar dataset, including the high-spatial $\times$
strong-temporal regime (\texttt{skel+gait}, EC1 $0.007$) and the
heavy-tailed \texttt{burr\_m5} (EC1 $0.004$).  KS--Flow tracks within a
factor of two and narrowly wins on \texttt{gait\_periodic}.
The conditional Gaussian oracle row (Table~\ref{tab:synth_ec1}) shows
residual MC noise: the oracle \emph{method} uses $n_\text{mc}{=}10$
random conditional samples per coalition (the same inference budget as
KernelSHAP imputers), and its EC1 of $0.06$--$0.12$ represents the
variance floor of this low-sample estimator.  Learned imputers
(KS--VAEAC, KS--Flow) substantially undercut the oracle method row
because each VAEAC sample approximates the conditional mean far more
accurately than a single Gaussian draw: effectively, the generator
``pre-conditions'' all relevant structure, making $M{=}1$ completion
nearly as good as $M{=}n_\text{mc}$.  Scaling to $n_\text{mc}{=}200$ is
predicted to reduce oracle EC1 by $\sqrt{10}\approx 3.16\times$,
converging toward the oracle's noiseless limit.  The learned imputers
already lie \emph{below} this limit at all tested datasets.
\textbf{(ii)} \emph{Off-manifold imputers cluster around a single error
level} (KS--Zero / KS--Mean / KS--Marginal $\in [0.08, 0.25]$), with no
clear winner; KS--Empirical (kernel-weighted donor sampling) is
competitive only on \texttt{gait\_periodic}.
\textbf{(iii)} \emph{Temporal-window players alone do not improve
attribution.}  KS--Temporal tracks KS--Zero closely; WindowSHAP performs
worst (EC1 $4$--$6\times$ the next-worst KernelSHAP variant).
WindowSHAP's failure is structural and not fixable by tuning window size:
Table~\ref{tab:winshap_diag} (Appendix) shows EC1 increasing monotonically
from $0.24$ ($w{=}2$) to $2.6$ ($w{=}16$) on \texttt{gauss\_k4}, and
from $0.67$ ($w{=}2$) to $1.6$ ($w{=}16$) on \texttt{skel+gait} ---
always $10$--$100\times$ above KS--VAEAC.  The root cause is that
WindowSHAP's sliding window-mean baseline assigns attribution mass to
boundary coalitions rather than interior frames, violating the
boundary-coalition constraint that KernelSHAP's weighted-least-squares
enforces, independent of $w$.
Temporal-aware coalition structure without manifold-aware imputation
offers no measurable benefit.
\textbf{(iv)} \emph{Gradient baselines (IG, DeepLIFT, SmoothGrad, LRP) are
not competitive on EC1}, with errors typically $5$--$50\times$ above
off-manifold KernelSHAP.  This is expected (gradient methods estimate
sensitivity, not cooperative-game value), but makes the oracle-anchored
cost of choosing a non-Shapley attribution explicit.

\paragraph{Player-set granularity.}
Comparing \texttt{gauss\_k4} ($K{=}4$) and \texttt{gauss\_k8} ($K{=}8$)
isolates the effect of player granularity: the on-manifold dominance is
unchanged, although $K{=}8$ doubles the number of coalitions ($2^{8}
{=}256$).  Table~\ref{tab:winners} (rightmost column) reports the best
EC1 across player-set granularities
($\mathcal{P}_{\mathrm{temp}}$, $\mathcal{P}_{\mathrm{joint}}$,
$\mathcal{P}_{\mathrm{cell}}$) for KS--VAEAC on the two high-spatial
regimes.  For \texttt{skeleton} (high spatial, AR(1) temporal),
$\mathcal{P}_{\mathrm{joint}}$ (17 spatial players) yields EC1 $\approx
0.010$--$0.020$ across classifiers, compared to $\mathcal{P}_{\mathrm{temp}}$
(4 temporal players) at $0.004$; the coarser temporal partition wins
because KernelSHAP's coalition-sampling budget covers 4 players far
better than 17, and VAEAC's generator already captures spatial structure
implicitly.  For \texttt{skel+gait} (high spatial, strong temporal), a
similar effect holds for $\mathcal{P}_{\mathrm{cell}}$ (68
joint$\times$window players).  In both regimes, on-manifold EC1 stays
below $0.007$ for $\mathcal{P}_{\mathrm{temp}}$, confirming that
on-manifold imputation benefit is robust to player-set choice and that
granularity should be matched to interpretability needs rather than
optimised for EC1.

\begin{table}[t]
\centering\small
\caption{Attribution error EC1 ($\downarrow$) on synthetic datasets, averaged over three classifier architectures (MLP, CNN, Transformer); mean $\pm$ std across classifiers.  EC1 is the mean absolute Shapley error \citep{aas2021explaining,olsen2022using} between estimated and oracle Shapley values per player.  Methods are grouped into KernelSHAP imputation variants (off- and on-manifold), temporal-SHAP baselines (KS--Temporal, WindowSHAP), and gradient-based attributions (Integrated Gradients, DeepLIFT, SmoothGrad, LRP).  \textbf{Bold} marks the best non-oracle method per dataset.  Oracle row shows the conditional Gaussian estimator's EC1 (Monte-Carlo noise floor of the reference oracle; lower is better).  Italic $M$-ablation rows show KS--VAEAC EC1 on \texttt{gauss\_k4} as $M{=}1$ completion samples used in the main rows; EC1 saturates by $M{=}5$.  Empty cells (--) indicate no result available.}
\label{tab:synth_ec1}
\begin{tabular}{lrrrrr}
\toprule
Method & \texttt{gauss\_k4} & \texttt{skeleton} & \texttt{gait} & \texttt{skel+gait} & \texttt{burr\_m5} \\
\midrule
KS--Zero & $0.101\!\pm\!0.055$ & $0.084\!\pm\!0.051$ & $0.143\!\pm\!0.122$ & $0.188\!\pm\!0.060$ & $0.087\!\pm\!0.048$ \\
KS--Mean & $0.102\!\pm\!0.055$ & $0.091\!\pm\!0.055$ & $0.139\!\pm\!0.116$ & $0.185\!\pm\!0.058$ & $0.087\!\pm\!0.048$ \\
KS--Marginal & $0.119\!\pm\!0.058$ & $0.114\!\pm\!0.072$ & $0.131\!\pm\!0.104$ & $0.164\!\pm\!0.046$ & $0.103\!\pm\!0.048$ \\
KS--Empirical & $0.122\!\pm\!0.044$ & $0.108\!\pm\!0.056$ & $0.057\!\pm\!0.047$ & $0.085\!\pm\!0.036$ & $0.112\!\pm\!0.035$ \\
KS--VAEAC & \textbf{$0.006\!\pm\!0.008$} & \textbf{$0.004\!\pm\!0.006$} & $0.005\!\pm\!0.006$ & \textbf{$0.007\!\pm\!0.008$} & \textbf{$0.004\!\pm\!0.005$} \\
KS--Flow & $0.010\!\pm\!0.014$ & $0.010\!\pm\!0.015$ & \textbf{$0.004\!\pm\!0.005$} & $0.007\!\pm\!0.009$ & $0.007\!\pm\!0.009$ \\
\midrule
KS--Temporal & $0.099\!\pm\!0.050$ & $0.083\!\pm\!0.051$ & $0.142\!\pm\!0.122$ & $0.187\!\pm\!0.059$ & $0.085\!\pm\!0.046$ \\
WindowSHAP & $0.465\!\pm\!0.213$ & $0.382\!\pm\!0.259$ & $0.544\!\pm\!0.620$ & $\mathit{1.037}^{\dagger}$ & $0.445\!\pm\!0.251$ \\
\midrule
IG & $1.217\!\pm\!0.820$ & $0.944\!\pm\!0.734$ & $1.344\!\pm\!1.846$ & $1.851\!\pm\!1.162$ & $0.977\!\pm\!0.658$ \\
DeepLIFT & $0.837\!\pm\!0.690$ & $0.622\!\pm\!0.603$ & $0.709\!\pm\!0.749$ & $1.419\!\pm\!1.110$ & $0.781\!\pm\!0.544$ \\
SmoothGrad & $2.351\!\pm\!1.589$ & $2.900\!\pm\!2.347$ & $4.713\!\pm\!3.453$ & $9.520\!\pm\!7.058$ & $2.775\!\pm\!2.011$ \\
LRP & $0.381\!\pm\!0.327$ & $0.405\!\pm\!0.445$ & $0.057\!\pm\!0.008$ & $0.251\!\pm\!0.226$ & $0.175\!\pm\!0.037$ \\
\midrule
Oracle$^{\ddagger}$ & $0.097\!\pm\!0.048$ & $\mathit{0.018}^{\dagger}$ & -- & $0.081\!\pm\!0.054$ & $0.062\!\pm\!0.044$ \\
\midrule
\multicolumn{6}{l}{\footnotesize\itshape $M$-ablation: KS--VAEAC completion samples on \texttt{gauss\_k4} (sensitivity saturates by $M{=}5$):} \\
\textit{KS--VAEAC ($M{=}1$)} & 0.0043 & -- & -- & -- & -- \\
\textit{KS--VAEAC ($M{=}5$)} & 0.0034 & -- & -- & -- & -- \\
\textit{KS--VAEAC ($M{=}20$)} & 0.0027 & -- & -- & -- & -- \\
\textit{KS--VAEAC ($M{=}50$)} & 0.0025 & -- & -- & -- & -- \\
\bottomrule
\end{tabular}
\vspace{0.5ex}\par\footnotesize $^{\dagger}$Single-architecture result (only one of MLP/CNN/Transformer completed).  Oracle \texttt{skeleton} (0.018, MLP only): the low value is an artefact of both the oracle method and oracle reference sharing the same approximate KernelSHAP WLS system with only $n_\text{coalitions}{=}64$ of $2^{17}$ possible coalitions; not directly comparable to low-$M$ datasets where exact enumeration is used.  WindowSHAP \texttt{skel+gait} (1.037, Transformer only): the default run errored due to a device-tensor issue on the other architectures; this entry uses window size $w{=}4$.
\vspace{0.5ex}\par\footnotesize $^{\ddagger}$Oracle samples $n_\text{mc}{=}10$ random conditional draws per coalition (same budget as the KernelSHAP imputers); EC1 $\approx 0.06$--$0.12$ is the MC noise floor of this estimator.  See \S\ref{sec:results-error}.
\end{table}


\subsection{The Spatial--Temporal Failure-Mode Matrix}
\label{sec:results-2x2}

\begin{table}[t]
\centering\small
\caption{Attribution error EC1 ($\downarrow$) of the best on-manifold and best off-manifold method in each spatiotemporal regime, with the \emph{on-manifold gap} $\Delta = \mathrm{EC1}_\mathrm{off} / \mathrm{EC1}_\mathrm{on}$.  The rightmost column reports the best KS--VAEAC EC1 across player-set granularities ($\mathcal{P}_\text{temp}$, $\mathcal{P}_\text{joint}$, $\mathcal{P}_\text{cell}$), with the winning granularity in parentheses.  For high-spatial regimes ($J{=}17$), $\mathcal{P}_\text{joint}$ (17 players) and $\mathcal{P}_\text{cell}$ (68 players) were compared against $\mathcal{P}_\text{temp}$ (4 players); for low-spatial datasets ($J{=}5$) only $\mathcal{P}_\text{temp}$ was evaluated, as the alternative granularities add negligible players.  On-manifold imputers dominate every regime by 1--2 orders of magnitude.}
\label{tab:winners}
\begin{tabular}{lllll}
\toprule
Regime & Best on-manifold & Best off-manifold & Gap $\Delta$ & Best player-set \\
\midrule
low spatial, AR(1) temporal & KS--VAEAC (0.004) & KS--Temporal (0.085) & $19.6\times$ & 0.0061 ($\mathcal{P}_\text{temp}$) \\
low spatial, strong temporal & KS--Flow (0.004) & KS--Marginal (0.131) & $31.4\times$ & 0.0049 ($\mathcal{P}_\text{temp}$) \\
high spatial, AR(1) temporal & KS--VAEAC (0.004) & KS--Temporal (0.070) & $15.7\times$ & 0.0044 ($\mathcal{P}_\text{temp}$) \\
high spatial, strong temporal & KS--VAEAC (0.007) & KS--Marginal (0.164) & $25.1\times$ & 0.0066 ($\mathcal{P}_\text{temp}$) \\
\bottomrule
\end{tabular}
\end{table}


Table~\ref{tab:winners} quantifies the on-manifold benefit per regime:
the best on-manifold method beats the best off-manifold KernelSHAP
variant by $15$--$31\times$ on every quadrant tested, including the
high-spatial $\times$ strong-temporal regime (\texttt{skel+gait},
$25\times$) that most closely resembles real motion data.  KS--VAEAC
dominates three of the four regimes; KS--Flow narrowly dominates the
low-spatial strongly-temporal periodic regime, consistent with its
continuous-time formulation more naturally capturing periodic structure.
The off-manifold winner is regime-dependent (KS--Temporal in the two
AR(1) regimes, KS--Marginal in the two strongly-temporal regimes), but
absolute EC1 remains an order of magnitude above the on-manifold methods
regardless of which off-manifold method is chosen.
Heavy-tailed Burr-marginal data (\texttt{burr\_m5}, \texttt{burr\_m10})
is reported in Table~\ref{tab:synth_ec1}: we retrained VAEAC and Flow
imputers from scratch on the Burr-marginal distribution using the same
transformer architecture ($80$ epochs; see \texttt{REPRODUCIBILITY.md}).
KS--VAEAC and KS--Flow both reach EC1 $\le 0.005$ on \texttt{burr\_m5}
and \texttt{burr\_m10} despite tail index $ck{=}4$, matching their
Gaussian-pillar performance and confirming that on-manifold dominance is
not an artefact of light tails.

\subsection{Faithfulness and Ranking Diagnostics}
\label{sec:results-faithfulness}

Attribution-error metrics require ground-truth Shapley values and therefore
cannot be evaluated on real data.  We complement EC1 with two metrics that
\emph{do} transfer: faithfulness correlation and PlayerAOPC.  Both are
computed from the same coalition value vector $v(S)$ that KernelSHAP
already evaluates, so the additional cost is negligible.  On synthetic
data, faithfulness is moderately rank-correlated with EC1
(Spearman $\rho \approx 0.6$, supplementary), establishing that these
metrics are reasonable proxies in the absence of an oracle, but only as
\emph{relative} comparisons \emph{within a single value-function family}.

We highlight one caveat that motivates the multi-metric protocol of
Sec.~\ref{sec:results-real}: faithfulness is computed by removing players
in the order suggested by their Shapley values.  When the imputer produces
\emph{out-of-distribution} completions, removal causes large prediction
shifts that are easy to predict from any reasonable attribution, inflating
faithfulness without improving attribution accuracy.  This pathology is
most visible on real data where ground-truth Shapley values are unavailable.
The italic rows at the bottom of Table~\ref{tab:synth_ec1} show KS--VAEAC
EC1 as $M$ is varied from $1$ to $50$ on \texttt{gauss\_k4}: EC1 plateaus
by $M{=}5$ (gain $<0.001$ beyond that), justifying $M{=}1$ as the
main-sweep default while confirming the result is not an artefact of
underdrawing the conditional.
